\documentclass[11pt]{article}
\usepackage[a4paper,margin=27mm]{geometry}
\usepackage{amsmath,amssymb,amsthm,mathrsfs,hyperref}
\newtheorem{theorem}{Theorem}
\newtheorem{lemma}[theorem]{Lemma}
\newtheorem{proposition}[theorem]{Proposition}
\newcommand{\PSH}{\operatorname{PSH}}
\newcommand{\EE}{\mathcal E^1_\theta}
\newcommand{\CC}{\mathcal C}
\title{Approximating every intrinsic target by actual outputs\protect\\
before proving variational attainment}
\author{Unreviewed conditional research, continuation 199}
\date{6 October 2026}
\begin{document}
\maketitle
\begin{abstract}
Under the explicit genuine scalar, output and uniform variation
foundations, every intrinsic finite-dual-energy target is a weak limit
of outputs of normalized model-envelope near maximizers with bounded
energy costs. This conclusion uses approximate stationarity before
any maximizer is known to exist. We first justify model approximation
of the objective using Radon regularity for directed continuous costs;
ordinary measurable monotone convergence cannot simply be applied
to arbitrary nets. Uniform model variation extends to every genuine
base with the same constants. If the separate ordered-gap tail package
holds, every fixed genuine finite-energy test also has converging
relative integrals along the output approximants. These are conditional
bridges, and neither variational attainment nor the full inverse
equation is claimed.
\end{abstract}

\section{Actual domain and the precise unproved foundations}
Retain every smooth integral projective $n$-fold, $n\geq1$, over every
algebraically closed trivially valued characteristic-zero field, and
every actual big rational bundle $L$. Put $K=X^{\mathrm{an}}$,
$\theta=c_1(L)$, $V=\operatorname{vol}(L)>0$, and let $v_0$ be the
generic trivial valuation. All supported semivaluations remain on $K$.
The genuine cone consists of actual functions obtained from one
all-point decreasing FS net on varying actual rational bundles with
classes tending to $\theta$ and a finite generic endpoint.
No abstract energy completion, countability or probability division
replaces these data.

The following are explicit conditional inputs. They have been proposed
in the unreviewed partial papers; their historical status supplies no
theorem grant.
\begin{enumerate}
\item The genuine cone is convex and decreasing closed, with actual
obstacle envelopes whenever there is a genuine minorant. Its normalized
reference $R=P_\theta(0)$ satisfies $R(v_0)=0$ and $u\leq R$ for
normalized $u$. The raw section energy is concave, monotone and
decreasing continuous, with $E(R)=0$ and $E(u+c)=E(u)+cV$.
Write $B(u)=-E(u)$ for normalized genuine $u\in\EE$.
Every such $u$ has nonpositive rational model majorants
$f_i\downarrow u$, $f_i(v_0)=0$, whose genuine envelopes
$u_i=P_\theta(f_i)$ decrease to $u$ with $E(u_i)\downarrow E(u)$.
For fixed $h\in C(K)$ their perturbed envelopes decrease to
$P_\theta(u+th)$ for each fixed $t$.
\item One actual Kodaira divisor $Z=E_0^{\mathrm{an}}$ has open
complement $U$. The reference and every bounded continuous-obstacle
envelope are finite continuous on $U$. There is a genuine normalized
$w_Z\in\EE$ with $w_Z=-\infty$ on $Z$.
For some dimension constant $C_n\geq1$, the genuine contraction obeys
\[
 B((1-r)R+ru)\leq C_n r^2 B(u),\qquad0<r<1.                 \tag{1}
\]
\item A target-independent map $M$ has positive original Radon
outputs of raw mass $V$ on all genuine $\EE$. At a model envelope
it equals its finite divisorial volume-derivative measure. For the
canonical majorants in item1, $M(u_i)\to M(u)$ weakly.
All outputs give $Z$ zero mass.
\item The proposed uniform finite-model Taylor estimate is valid.
For a fixed rational model direction $h$, a cost bound $B$, and
$c>B/V$, there are $K_c<\infty$ and $t_c>0$, independent of the
nonpositive model obstacle and its refinement, such that
\[
 \left|\mathscr E(f+th)-\mathscr E(f)
                -t\int h\,dM(P_\theta(f))\right|
 \leq K_ct^2/2+|t|A_h(B)/c,\qquad |t|<t_c,                 \tag{2}
\]
when $f(v_0)=0$ and $B(P_\theta(f))\leq B$.
Here $\mathscr E(f)=E(P_\theta(f))$ and one may use
$A_h(B)=\|h\|_\infty(2B+2V\|h\|_\infty)$.
Rational model functions are uniformly dense in the original $C(K)$.
\end{enumerate}
Items1--4 are the antecedents of the approximation theorem below.
In particular (2) is stronger than pointwise volume differentiability.
Its proposed proof uses deficient positive measures with tail mass
$\leq B/c$, and fixed signed Cartier tests paired through two fixed
nef tests on actual big smooth algebraic determinations. Those proofs
remain to be independently audited. Neither an ample unbounded
estimate nor an inverse equation is assumed here.

Let $\nu$ be any positive original Radon target of mass $V$ with
literal reference-pole exclusion, all relative finite-energy costs
finite, and finite intrinsic dual energy:
\[
 \begin{gathered}
 \nu\{R=-\infty\}=0,\quad
 I_\nu(u)=\int(R-u)\,d\nu<\infty\ (u(v_0)=0),\\
 D_\nu=\sup_{u\in\EE,\ u(v_0)=0}\{E(u)+I_\nu(u)\}<\infty.
 \end{gathered}                                             \tag{3}
\]
Relative values are actual differences where $R$ is finite.
The existence of $w_Z$ implies $\nu(Z)=0$: its cost is infinite
on $Z$ outside reference poles, and the latter are already excluded.
No absolute integrability of $R$ is needed. The earlier capped-section
rank construction proposes $w_Z$ in precisely this genuine space.

\section{Directed continuous costs and near maximizers}
\begin{lemma}[Radon convergence for the actual directed costs]\label{radon}
Let $c_i$ be a directed increasing net of nonnegative continuous
functions on $U$, and let $\nu$ be finite Radon on $K$, $\nu(Z)=0$.
Then $\int_U\sup_i c_i\,d\nu=\sup_i\int_U c_i\,d\nu$,
with $+\infty$ permitted.
\end{lemma}
\begin{proof}
For $a>0$, $\{\sup_i c_i>a\}=\bigcup_i\{c_i>a\}$ is open in $U$.
Every compact subset of this union has a finite subcover by these
open superlevels, and directedness gives one index dominating that
finite set of indices. Radon inner regularity therefore gives
\[
 \nu\{\sup_i c_i>a\}=\sup_i\nu\{c_i>a\}.
\]
For any finite collection of thresholds the corresponding masses can
be approached simultaneously by one index. Apply this fact to the
finite lower step sums for the layer-cake integral of
$\min(\sup_i c_i,N)$. Then refine the threshold mesh and send
$N\to\infty$. It proves the nontrivial inequality; the other follows
from $c_i\leq\sup_jc_j$. This argument uses Radon regularity and
continuity, rather than measurable monotone convergence for arbitrary
uncountable nets.
\end{proof}

For the actual majorants of any normalized $u\in\EE$, the costs
$R-u_i$ are nonnegative bounded continuous on $U$, increase there
to $R-u$, and have finite $\nu$ integral. Lemma~\ref{radon} and
decreasing energy continuity show
\[
 E(u_i)+I_\nu(u_i)\longrightarrow E(u)+I_\nu(u).             \tag{4}
\]
Thus taking the supremum only over actual normalized model envelopes
still gives exactly $D_\nu$. For each $k\geq1$ choose one such
$u_k$ with objective at least $D_\nu-1/k$.
This choice is a sequence of near maximizers of a real supremum;
it does not impose metrizability on $K$ or replace the defining nets.

The costs of these near maximizers are uniformly bounded. Indeed
$D_\nu\geq0$ because $R$ has objective zero. By (1) and linearity
of the relative integral on a genuine convex combination,
\[
 r I_\nu(u)-C_nr^2B(u)\leq D_\nu.
\]
At $r=1/(2C_n)$ this gives
\[
 I_\nu(u)\leq B(u)/2+2C_nD_\nu,\qquad
 E(u)+I_\nu(u)\leq-B(u)/2+2C_nD_\nu.                      \tag{5}
\]
Consequently every chosen $u_k$ has
$B(u_k)\leq4C_nD_\nu+2=:B_\nu$.
This is coercivity derived from the explicit contraction input;
bounded cost alone supplies no variational attainment.

\section{The same uniform variation estimate at every genuine base}
\begin{proposition}\label{variation}
Under items1--4, (2) extends with the same $B,c,h,K_c,t_c,A_h(B)$
to every normalized genuine $u\in\EE$ with $B(u)\leq B$:
\[
 \left|E(P_\theta(u+th))-E(u)-t\int h\,dM(u)\right|
 \leq K_ct^2/2+|t|A_h(B)/c.                                \tag{6}
\]
\end{proposition}
\begin{proof}
Take its actual nonpositive model majorants $f_i$ from item1.
Their envelope costs are at most $B$, so (2) holds uniformly for
all $i$ with the same constants. The perturbed envelopes lie between
$u-|t|\|h\|_\infty$ and $u+|t|\|h\|_\infty$ at their endpoint,
and hence have finite genuine energy. At fixed $t$, their energies
decrease to $E(P_\theta(u+th))$ by item1. The unperturbed energies
decrease to $E(u)$, and $\int h\,dM(u_i)\to\int h\,dM(u)$
by item3. Taking this limit in (2) proves (6).
There is no passage of a derivative through the net: the entire
uniform finite increment estimate passes at fixed $t$ first.
\end{proof}

\section{Approximate stationarity gives the target before attainment}
\begin{theorem}\label{approx}
Under items1--4, every target (3) is a weak limit of actual outputs
$M(u_k)$ of normalized model-envelope near maximizers. Their costs
are bounded by $B_\nu$, and no actual maximizing potential is
assumed to exist.
\end{theorem}
\begin{proof}
Fix one rational model $h$. Write $\epsilon_k=1/k$ and
$z_{k,t}=P_\theta(u_k+th)$. This is a genuine finite-energy
potential by the bounds in Proposition~\ref{variation}. Its generic
value is finite; normalize it when applying (3). Raw energy and mass
make the objective invariant under that constant shift.
The obstacle inequality $z_{k,t}\leq u_k+th$ gives, on $U$,
\[
 I_\nu(z_{k,t})\geq I_\nu(u_k)-t\int h\,d\nu.
\]
All terms are finite by (3), and the near-maximizer inequality yields
\[
 E(z_{k,t})-E(u_k)-t\int h\,d\nu\leq\epsilon_k.             \tag{7}
\]
It is essential that the perturbed obstacle here is $u_k+th$,
rather than an arbitrary larger model majorant: the latter would
introduce an uncontrolled integral against $\nu$.

Use (6) at the cost bound $B_\nu$ for both signs of $t$, with
fixed $c>B_\nu/V$. For positive $t<t_c$ it gives
\[
 \left|\int h\,d(M(u_k)-\nu)\right|
 \leq\epsilon_k/t+K_ct/2+A_h(B_\nu)/c.                    \tag{8}
\]
First send $k\to\infty$ at fixed $c,t$, then $t\downarrow0$
for this fixed $c$, and finally $c\to\infty$.
Both one-sided integral limits coincide with $\int h\,d\nu$.
Original uniform model density and the common raw mass $V$ give
the same limit for every $h\in C(K)$. This is weak convergence of
the whole chosen sequence on the original compact carrier.
No compact-cluster maximizer or fixed-target strong integration
has been used.
\end{proof}

\section{Optional fixed-test integration consequence}
The following needs the additional, explicitly conditional ordered-gap
tail package proposed in the completed partial paper. For every
normalized $v\in\EE$, set $b_v=R-v$ on $U$ and zero on $Z$.
For every cost bound $B$ and $\eta>0$ assume that an actual
$a_\eta\in C(K)$ satisfies
\[
 0\leq a_\eta\leq b_v,\quad a_\eta|_Z=0,\qquad
 \sup_{B(u)\leq B}\int(b_v-a_\eta)\,dM(u)<\eta.             \tag{9}
\]
The proposed construction first fixes a model majorant of $v$ using
its ordered energy gap, then fixes a continuous Kodaira collar using
the uniform algebraic-cap estimate. This fixes both choices before
any output index. Formula (9) is stronger than a bounded first moment.

\begin{proposition}\label{integration}
If (9) holds, the output sequence in Theorem~\ref{approx} satisfies,
for every fixed genuine $v\in\EE$ and $h\in C(K)$,
\[
 \int h(R-v)\,dM(u_k)\longrightarrow\int h(R-v)\,d\nu.     \tag{10}
\]
The family consisting of these outputs and $\nu$ has uniform
integrable tails for this fixed relative test.
\end{proposition}
\begin{proof}
On $U$, $b_v$ is nonnegative lsc, since $R$ is continuous and $v$
is usc. Assigning zero on $Z$ gives a global nonnegative lsc
function. Thus $b_v-a_\eta$ is nonnegative lsc on $K$.
Weak Radon Portmanteau, valid on compact Hausdorff $K$, and (9) give
\[
 \int(b_v-a_\eta)\,d\nu
       \leq\liminf_k\int(b_v-a_\eta)\,dM(u_k)\leq\eta.
\]
The target is algebraically null, so this is the actual relative
integral in (3), without hiding values on reference poles.
The integral of the fixed continuous $h a_\eta$ passes weakly.
Both error terms are at most $\|h\|_\infty\eta$.
Let $\eta\downarrow0$ to prove (10).
If $C=\|a_\eta\|_\infty$, then
$b_v\mathbf1_{\{b_v>2C\}}\leq2|b_v-a_\eta|$.
Choose $\eta$ first and then this fixed threshold to prove uniform
tails over the stated family. Shifts of $v$ use the raw mass law.
\end{proof}

Theorem~\ref{approx} places every intrinsic target in a weak closure
of a bounded-cost actual output family. Proposition~\ref{integration}
then gives convergence for each fixed genuine test. Neither assertion
controls a simultaneously changing test $u_k$ against $\nu$, or makes
$E(u_k)+I_\nu(u_k)$ attain its supremum at a genuine cluster.
Strict concavity, a genuine strong-completeness argument and exclusion
of all energy/integral jumps are still required. Identification of
all mixed unbounded moving products, strict Borel comparison and
uniqueness of every solution also remain open.
This conditional paper does not close original P16.5 or change the
accepted-paper count. No active campaign output was read or hot
feedback sent.

\begin{thebibliography}{9}\small
\bibitem{scalar} Whole 14-page unreviewed genuine-cone, section-volume
energy, contraction and outer-output foundation paper, 6 October 2026,
especially the uniform finite-model variation estimate (11.7).
\bibitem{bounded} Bounded-obstacle full-section density and general
bounded relative weighted-limit draft, continuation198, 6 October 2026.
\bibitem{gap} Whole completed 12-page unreviewed partial paper,
\emph{Bounded-obstacle section approximation and ordered energy-gap tails},
6 October 2026; its output package F is an explicit unproved antecedent.
\end{thebibliography}
\end{document}
